\documentclass[11pt]{article}

\usepackage[margin=1in]{geometry}
\usepackage{amsmath,amssymb,amsthm,mathtools,bm}
\usepackage{mathrsfs}
\usepackage{hyperref}
\usepackage{enumitem}
\usepackage{microtype}

\newtheorem{theorem}{Theorem}[section]
\newtheorem{proposition}[theorem]{Proposition}
\newtheorem{lemma}[theorem]{Lemma}
\newtheorem{corollary}[theorem]{Corollary}
\theoremstyle{definition}
\newtheorem{definition}[theorem]{Definition}
\theoremstyle{remark}
\newtheorem{remark}[theorem]{Remark}

\newcommand{\R}{\mathbb{R}}
\newcommand{\E}{\mathbb{E}}
\newcommand{\Lip}{\operatorname{Lip}}
\newcommand{\Id}{\operatorname{Id}}
\newcommand{\dd}{\,\mathrm{d}}
\newcommand{\norm}[1]{\left\lVert #1\right\rVert}
\newcommand{\abs}[1]{\left\lvert #1\right\rvert}
\newcommand{\avgint}{\mathop{\rlap{\raisebox{.15ex}{--}}\!\int}\nolimits}
\newcommand{\hsemig}[2]{\mathcal{G}_{#1}\left[{#2}\right]}
\newcommand{\V}{\mathcal{V}}
\newcommand{\Zop}{\mathcal{Z}}
\newcommand{\X}{\mathscr{X}}
\newcommand{\Hs}{\mathcal{H}}
\newcommand{\JF}{J_{\sigma,\rho}}
\newcommand{\J}{J_{\sigma}}
\newcommand{\Trho}{T_{\rho}}
\newcommand{\Srho}{S_{\rho}}
\newcommand{\Lrho}{L_{\rho}}
\newcommand{\Kpar}[1]{K_{\lambda,\rho,#1}}

\hypersetup{colorlinks=true,linkcolor=blue,citecolor=blue,urlcolor=blue}

\title{Chapter 5 Notes: Local SPDE Estimates}
\author{Wells Guan}
\date{}

\begin{document}
\maketitle

\section{Local SPDE Estimates}

This chapter records the local estimates from Chapter 5 of
\cite{dunlap2026renormalizationflow2dnonlinear} in the notation used in our bounded-nonlinearity project.
We consider the post-smoothed stochastic heat equation
\begin{equation}\label{eq:ch5-spde}
 v_t^{\rho}(x)
 = \hsemig{t-s}{v_s^{\rho}}(x)
 + \gamma_{\rho}\int_s^t
 \hsemig{t+\rho-r}{\sigma(v_r^{\rho})\dd W_r}(x),
\end{equation}
where
\[
\gamma_{\rho}=\left(\frac{4\pi}{L(1/\rho)}\right)^{1/2},
\qquad
L(\tau)=\log(1+\tau).
\]
As in the paper, define
\begin{equation}\label{eq:ch5-scales}
\Lrho(\tau)=\frac{L(\tau)}{L(1/\rho)},
\qquad
\Srho(\tau)=\Lrho(\tau/\rho)
=\frac{L(\tau/\rho)}{L(1/\rho)},
\end{equation}
and
\begin{equation}\label{eq:ch5-Trho}
\Trho(q)=\Srho^{-1}(q)
=\rho\left[(1+1/\rho)^q-1\right].
\end{equation}
For an upper bound \(\lambda\ge \Lip(\sigma)\), introduce the short time scale
\begin{equation}\label{eq:ch5-Tbar}
\overline T_{\rho}(\lambda)
:=\Trho\!\left(\lambda^{-2}-2L(1/\rho)^{-1}\right)
\end{equation}
and the amplification factor
\begin{equation}\label{eq:ch5-K}
K_{\lambda,\rho,t}
:=\left(1-\lambda^2\left[\Srho(t)+2L(1/\rho)^{-1}\right]\right)^{-1},
\qquad 0\le t<\overline T_{\rho}(\lambda).
\end{equation}
Thus the estimates below are local in scale: roughly speaking, they work while
\[
\Srho(t)<\lambda^{-2},
\qquad\text{equivalently}\qquad
 t\lesssim \rho^{1-\lambda^{-2}}.
\]
The factor \(K_{\lambda,\rho,t}\) remains bounded away from the endpoint but diverges as the endpoint is approached.

\subsection{Turning off the noise and flattening}

The first operation is to insert a quiet interval, during which the stochastic forcing is switched off and only the deterministic heat flow is run. The second operation replaces the field after the quiet interval by a spatially constant field.

\subsubsection{Turning off the noise field}

For \(\tau_0\le \tau_1<\tau_2\), compare
\[
\V^{\sigma,\rho}_{\tau_0,\tau_2}v_{\tau_0}
\qquad\text{and}\qquad
\V^{\sigma,\rho}_{\tau_1,\tau_2}
\hsemig{\tau_1-\tau_0}{v_{\tau_0}}.
\]
The second evolution has no noise on \([\tau_0,\tau_1]\).

\begin{proposition}[Turning off the noise]\label{prop:ch5-turnoff}
Assume \(\Lip(\sigma)\le\lambda\) and
\[
\tau_2-\tau_1<\overline T_{\rho}(\lambda).
\]
Then the error produced by suppressing the noise on \([\tau_0,\tau_1]\) is bounded by
\begin{equation}\label{eq:ch5-turnoff-bound}
\begin{aligned}
&\norm{\V_{\tau_0,\tau_2}v_{\tau_0}
-\V_{\tau_1,\tau_2}\hsemig{\tau_1-\tau_0}{v_{\tau_0}}}_2^2
\\
&\quad\le
K_{\lambda,\rho,\tau_2-\tau_1}\,
\Xi_{2;\tau_0,\tau_2}[v_{\tau_0}]^2
\left[
\Lrho\!\left(\frac{\tau_1-\tau_0}{\tau_2-\tau_1+\rho}\right)
+\frac{L(2)K_{\lambda,\rho,\tau_2-\tau_1}}{L(1/\rho)}
\right].
\end{aligned}
\end{equation}
\end{proposition}

\begin{proof}
Set \(\tau_0=0\) for simplicity and let
\[
v_t=\V_{0,t}v_0,
\qquad
\widetilde v_t=\V_{\tau_1,t}\hsemig{\tau_1}{v_0}.
\]
The mild equations give
\[
 v_t-\widetilde v_t
 =\gamma_\rho\int_0^{\tau_1}
 \hsemig{t+\rho-r}{\sigma(v_r)\dd W_r}
 +\gamma_\rho\int_{\tau_1}^t
 \hsemig{t+\rho-r}{[\sigma(v_r)-\sigma(\widetilde v_r)]\dd W_r}.
\]
Using the It\^o isometry and
\[
4\pi\int_{\R^2}G_s(y)^2\dd y=\frac1s,
\]
we obtain a Volterra inequality of the form
\begin{equation}\label{eq:ch5-volterra1}
f(\tau)
\le A L\!\left(\frac{\zeta}{\tau+\rho}\right)
+\frac{\lambda^2}{L(1/\rho)}
\int_0^\tau\frac{f(r)}{\tau+\rho-r}\dd r.
\end{equation}
The next lemma closes this estimate. Applying it with the parameters coming from the first stochastic integral yields \eqref{eq:ch5-turnoff-bound}.
\end{proof}

\begin{lemma}[Logarithmic Volterra estimate]\label{lem:ch5-volterra}
Suppose \(f:[0,T]\to[0,\infty)\) is bounded and satisfies \eqref{eq:ch5-volterra1}. Then for
\(t<T\wedge\overline T_{\rho}(\lambda)\),
\begin{equation}\label{eq:ch5-volterra-result}
f(t)
\le A K_{\lambda,\rho,t}
\left[
L\!\left(\frac{\zeta}{t+\rho}\right)
+L(2)K_{\lambda,\rho,t}
\right].
\end{equation}
\end{lemma}

\begin{proof}
Use the ansatz
\[
f(t)\le B_1L\!\left(\frac{\zeta}{t+\rho}\right)+B_2.
\]
Split the Volterra integral into \([0,t/2]\) and \([t/2,t]\). The logarithmic kernel obeys
\[
\int_0^t
\frac{L(\zeta/(s+\rho))}{t+\rho-s}\dd s
\le
L\!\left(\frac{\zeta}{t+\rho}\right)
+L(2)[L(t/\rho)+2].
\]
Writing
\[
\alpha=\lambda^2\left(\Srho(t)+2L(1/\rho)^{-1}\right)<1,
\]
the coefficients satisfy a geometric recursion, and summing it gives the factor
\[
(1-\alpha)^{-1}=K_{\lambda,\rho,t}.
\]
\end{proof}

\subsubsection{Flattening}

For \(X\in\R^2\), define
\begin{equation}\label{eq:ch5-Z}
(\Zop_Xv)(x)=v(X),\qquad x\in\R^2.
\end{equation}
Thus \(\Zop_X\) replaces a field by its value at \(X\).

\begin{proposition}[Flattening estimate]\label{prop:ch5-flatten}
Assume \(\Lip(\sigma)\le\lambda\) and
\(\tau_2-\tau_1<\overline T_{\rho}(\lambda)\). Then for every \(x,X\in\R^2\),
\begin{align}\label{eq:ch5-flatten-bound}
&\E\abs{
\left(\V_{\tau_1,\tau_2}\hsemig{\tau_1-\tau_0}{v_{\tau_0}}
-\V_{\tau_1,\tau_2}\Zop_X\hsemig{\tau_1-\tau_0}{v_{\tau_0}}\right)(x)}^2
\\
&\quad\le
K_{\lambda,\rho,\tau_2-\tau_1}\norm{v_{\tau_0}}_2^2
\left[
\frac{\tau_2-\tau_1+\frac12|x-X|^2}{\tau_1-\tau_0}
+\frac{\lambda^2(\tau_2-\tau_1)K_{\lambda,\rho,\tau_2-\tau_1}}
{2(\tau_1-\tau_0)L(1/\rho)}
\right].
\end{align}
\end{proposition}

\begin{proof}
The deterministic difference is controlled by the heat kernel estimate
\[
\norm{G_{\tau_1+\tau}(x-\cdot)-G_{\tau_1}(X-\cdot)}_{L^1}^2
\lesssim \frac{\tau+|x-X|^2/2}{\tau_1}.
\]
The stochastic part yields
\[
f(\tau,x)
\le A_1+A_2|x-X|^2
+\frac{\lambda^2}{L(1/\rho)}
\int_0^\tau\int_{\R^2}
\frac{G_{(\tau+\rho-r)/2}(x-y)}{\tau+\rho-r}f(r,y)\dd y\dd r.
\]
The quadratic ansatz
\[
f(\tau,x)\le B_1+B_2|x-X|^2
\]
closes because in two dimensions
\begin{equation}\label{eq:ch5-quadratic-heat}
\int_{\R^2}G_{s/2}(x-y)|y-X|^2\dd y
=|x-X|^2+s.
\end{equation}
A geometric iteration again produces \(K_{\lambda,\rho,T}\), giving \eqref{eq:ch5-flatten-bound}.
\end{proof}

\begin{proposition}[Turning off and flattening together]\label{prop:ch5-combined}
Under the preceding hypotheses, the solution at time \(\tau_2\) can be compared with the solution which first runs the heat equation on \([\tau_0,\tau_1]\), is flattened at a point \(X\), and then restarts the SPDE. The total error is bounded by the sum of the turning-off error from Proposition~\ref{prop:ch5-turnoff} and the flattening error from Proposition~\ref{prop:ch5-flatten}.
\end{proposition}

\begin{proof}
Insert the intermediate field
\(
\V_{\tau_1,\tau_2}\hsemig{\tau_1-\tau_0}{v_{\tau_0}}
\)
and apply the triangle inequality in \(L^2(\Omega)\). This is the form used later in the proof of Theorem~\ref{thm:ch5-main-local}.
\end{proof}

\subsection{Local-in-space dependence on the noise}

The next input is that the value of the solution at a point depends only weakly on noise far away from that point. For a Borel set \(B\subset\R^2\), write \(\V^{B}_{s,t}\) for the SPDE evolution driven by the truncated noise \(\mathbf{1}_B\dd W\).

\begin{proposition}[Localization of the noise]\label{prop:ch5-localnoise}
There exists an absolute constant \(C\) such that the following holds. Suppose \(\Lip(\sigma)\le\lambda\), \(s<t\), and \(t-s<\overline T_{\rho}(\lambda)\). Let \(R\subset\R^2\) be a square and \(x\in R\). Then
\begin{equation}\label{eq:ch5-localnoise-bound}
\E\abs{\V_{s,t}a(x)-\V^R_{s,t}a(x)}^2
\le
C\,\Xi_{2;s,t}[a]^2\,
K_{\lambda,\rho,t-s}
\exp\!\left(-c\frac{\operatorname{dist}(x,R^c)^2}{t-s+\rho}\right),
\end{equation}
up to the harmless logarithmic factors recorded in the paper.
\end{proposition}

\begin{proof}
Write \(R^c\) as the union of four half-spaces. For each half-space, compare the full solution with the solution driven only by the noise on the complementary side. The It\^o isometry produces a Gaussian tail factor, while the feedback term is again handled by the short-time Volterra estimate. Summing the four contributions proves the square localization estimate.
\end{proof}

\subsection{The approximate root decoupling function: regularity}

\begin{definition}[Approximate root decoupling function]\label{def:ch5-Jrho}
For \(q\ge0\) and \(a\in\R^m\), define
\begin{equation}\label{eq:ch5-Jrho}
J_{\sigma,\rho}(q,a)
:=\left(\E\sigma^2\!\left(\V_{0,\Trho(q)}a(x)\right)\right)^{1/2}.
\end{equation}
By spatial stationarity, the right-hand side does not depend on \(x\).
\end{definition}

\begin{proposition}[Spatial Lipschitz estimate]\label{prop:ch5-Jrho-space}
If \(\Lip(\sigma)\le\lambda\), then for \(0<q<\lambda^{-2}\),
\begin{equation}\label{eq:ch5-Jrho-space}
\abs{J_{\sigma,\rho}(q,a_2)-J_{\sigma,\rho}(q,a_1)}_F
\le
\frac{|a_2-a_1|}{(\lambda^{-2}-q)^{1/2}}.
\end{equation}
\end{proposition}

\begin{proof}
Proposition A.2 of the paper reduces the difference of matrix square roots to
\[
\lambda^2\norm{\V_{\Trho(q)}a_2-\V_{\Trho(q)}a_1}_2^2.
\]
The difference of the two solutions obeys an It\^o-isometry inequality. Taking the supremum in time yields
\[
A\le |a_2-a_1|^2+\lambda^2qA,
\]
so
\[
A\le\frac{|a_2-a_1|^2}{1-\lambda^2q}.
\]
Substitution gives \eqref{eq:ch5-Jrho-space}.
\end{proof}

\begin{proposition}[Time regularity]\label{prop:ch5-Jrho-time}
Suppose \(q_1\wedge q_2<\lambda^{-2}-2L(1/\rho)^{-1}\). Then
\begin{align}\label{eq:ch5-Jrho-time}
&|J_{\sigma,\rho}(q_2,a)-J_{\sigma,\rho}(q_1,a)|_F
\\
&\quad\le
\lambda K_{\lambda,\rho,\Trho(q_1\wedge q_2)}^{1/2}
\Xi_{2;|\Trho(q_2)-\Trho(q_1)|}[a]
\left[
|q_2-q_1|^{1/2}
+\frac{L(2)^{1/2}K_{\lambda,\rho,\Trho(q_1\wedge q_2)}^{1/2}}
{L(1/\rho)^{1/2}}
\right].
\end{align}
\end{proposition}

\begin{proof}
Use Proposition~\ref{prop:ch5-turnoff} to compare the evolution over \(\Trho(q_2)\) with the evolution over \(\Trho(q_1)\), shifted to the terminal time. The identity
\[
\Lrho\!\left(\frac{\Trho(q_2)-\Trho(q_1)}{\Trho(q_1)+\rho}\right)=q_2-q_1
\]
converts the physical-time estimate into the scale-variable estimate.
\end{proof}

\begin{corollary}[Joint regularity]\label{cor:ch5-Jrho-joint}
Combining Propositions~\ref{prop:ch5-Jrho-space} and \ref{prop:ch5-Jrho-time} gives a bound for
\(|J_{\sigma,\rho}(q_2,a_2)-J_{\sigma,\rho}(q_1,a_1)|_F\)
by the sum of the time increment in \eqref{eq:ch5-Jrho-time} and a spatial term proportional to \(|a_2-a_1|\).
\end{corollary}

\subsection{Concentration of averages of the \texorpdfstring{$\sigma^2$}{sigma2} field}

The next step converts spatial averaging into an approximation of expectation. The central mechanism is localization: after the noise is truncated to well-separated boxes, the corresponding random fields become independent.

\begin{proposition}[Concentration over a square]\label{prop:ch5-concentration}
Let \(\ell\in(2,4]\). There is \(C=C(\ell,m)\) such that, for constant initial data \(a\), a square \(R\subset\R^2\) of side length \(\xi\), and
\[
t-s<\xi^2\wedge\overline T_{\rho}(\lambda),
\]
one has
\begin{align}\label{eq:ch5-concentration}
&\E\left|
\left(\avgint_R\sigma^2(\V_{s,t}a(x))\dd x\right)^{1/2}
-J_{\sigma,\rho}(\Srho(t-s),a)
\right|_F^2
\\
&\quad\le
C\,\Xi_{\ell;s,t}[a]^2
\left[
\left(1\wedge\log\frac1{\lambda^2\Srho(t-s)}\right)
\frac{\log\left(\xi[(t-s)\vee\rho]^{-1/2}\right)}
{\xi[(t-s)\vee\rho]^{-1/2}}
\right]^{2-4/\ell}.
\end{align}
\end{proposition}

\begin{proof}
Partition \(R\) into \(N^2\) smaller boxes. Enlarge each small box slightly and split the index set into finitely many color classes so that enlarged boxes in the same class are disjoint. Replace the full solution on each small box by a localized solution driven only by noise in the corresponding enlarged box. Proposition~\ref{prop:ch5-localnoise} controls this replacement. Within one color class, the localized random variables are independent, so an \(L^{\ell/2}\) estimate for sums yields a gain of order \(N^{-(2-4/\ell)}\). Optimize \(N\) against the Gaussian localization error to obtain \eqref{eq:ch5-concentration}.
\end{proof}

\begin{remark}
The exponent \(2-4/\ell\) is positive precisely because \(\ell>2\). If \(\Srho(t-s)\le c<\lambda^{-2}\), then the factor involving
\(\log(1/(\lambda^2\Srho(t-s)))\) is uniformly bounded.
\end{remark}

\begin{proposition}[Crude short-time bound]\label{prop:ch5-concentration-crude}
For arbitrary Borel \(R\subset\R^2\),
\begin{equation}\label{eq:ch5-crude}
\E\left|
\left(\avgint_R\sigma^2(\V_{s,t}a(x))\dd x\right)^{1/2}
-\sigma(a)
\right|_F^2
\lesssim \lambda^2\Xi_{2;s,t}[a]^2\Srho(t-s).
\end{equation}
This estimate is useful when \(t-s\lesssim\rho\), where the sharper concentration estimate is unnecessary.
\end{proposition}

\subsection{Approximating averages of \texorpdfstring{$\sigma^2$}{sigma2} by averages of the solution}

For \(\zeta>0\) and \(z\in\R^2\), let \(B_{\zeta,k}=\zeta(k+[0,1)^2)\) and define the box-average operator
\begin{equation}\label{eq:ch5-Sbox}
S_{\zeta,z}f(x)
:=\frac1{\zeta^2}\int_{B_{\zeta,k_{\zeta,z}(x)+z}}f(y)\dd y,
\end{equation}
where \(k_{\zeta,z}(x)\in\mathbb Z^2\) is the unique index whose box contains \(x\).
For \(\ell>2\), define
\begin{equation}\label{eq:ch5-kappa}
\kappa_{\ell,\rho}
:=L(1/\rho)^{-1/(1-2/\ell)}.
\end{equation}

\begin{theorem}[Local approximation of the \texorpdfstring{$\sigma^2$}{sigma2} field]\label{thm:ch5-main-local}
Fix \(M,\beta,\lambda>0\), \(\ell\in(2,4]\), choose \(\overline\beta>\eta(\beta,\ell)\), and choose \(\overline q\in(0,\lambda^{-2})\). Then there is
\[
C=C(M,\beta,\lambda,\ell,\overline\beta,\overline q)
\]
such that for sufficiently small \(\rho\), if
\[
\sigma\in\Sigma(M,\beta),\qquad \Lip(\sigma)\le\lambda,
\]
and \(v\) solves \eqref{eq:ch5-spde} with \(v_0\in\X_0^\ell\), then for admissible \(\xi,t,x,z\) and \(\iota_1,\iota_2\in[0,1]\),
\begin{align}\label{eq:ch5-theorem514}
&\E\left|
\big(S_{\xi,z}[\sigma^2\circ v_t](x)\big)^{1/2}
-J_{\sigma,\rho}\!\left(
\Srho(L(1/\rho)^{\iota_1}\xi^2),
\hsemig{L(1/\rho)^{\iota_2}\xi^2}{v_t}(x)
\right)
\right|_F^2
\\
&\qquad\le
C\,\left\langle\norm{v_0}_\ell\right\rangle^2
\frac{L(L(1/\rho))}{L(1/\rho)}.
\end{align}
The admissible range is the one in Theorem 5.14 of the paper; in particular \(\xi\) is below the local scale determined by \(\overline q\), and \(t\) is separated from the initial time by the quiet/flattening buffer of order
\([1+L(1/\rho)]\kappa_{\ell,\rho}^{-1}\xi^2+\kappa_{\ell,\rho}\xi^2\).
\end{theorem}

\begin{proof}
The proof is the main synthesis of Sections 5.1--5.4. Choose the three times
\[
\tau_0=t-(\kappa_{\ell,\rho}+\kappa_{\ell,\rho}^{-1})\xi^2,
\qquad
\tau_1=\tau_0+\kappa_{\ell,\rho}^{-1}\xi^2,
\qquad
\tau_2=\tau_1+\kappa_{\ell,\rho}\xi^2=t.
\]
Let \(X\) be the center of the averaging square.

First, apply the combined turning-off/flattening estimate. The resulting error is of the schematic form
\begin{equation}\label{eq:ch5-E1}
E_1
\lesssim
K_{\lambda,\rho,\xi^2}\Xi_{2;\tau_0,\tau_1}[v_{\tau_0}]^2
\left[
\Lrho(\kappa_{\ell,\rho}^{-2})
+\kappa_{\ell,\rho}\big((1+\lambda^2\delta)\kappa_{\ell,\rho}+1\big)
+\frac{K_{\lambda,\rho,\xi^2}}{L(1/\rho)}
\right],
\end{equation}
where
\[
\delta=K_{\lambda,\rho,\kappa_{\ell,\rho}\xi^2}L(1/\rho)^{-1}.
\]
Since \(\overline q<\lambda^{-2}\), the \(K\)-factors are uniformly bounded. Moreover,
\[
\Lrho(\kappa_{\ell,\rho}^{-2})
\lesssim \frac{L(L(1/\rho))}{L(1/\rho)},
\]
and the other terms are smaller. Hence
\begin{equation}\label{eq:ch5-E1final}
E_1\lesssim
\Xi_{2;\tau_0,\tau_1}[v_{\tau_0}]^2
\frac{L(L(1/\rho))}{L(1/\rho)}.
\end{equation}

Second, apply Proposition~\ref{prop:ch5-concentration} if
\(\kappa_{\ell,\rho}\xi^2\ge\rho\), and Proposition~\ref{prop:ch5-concentration-crude} otherwise. In either case,
\begin{equation}\label{eq:ch5-E2final}
E_2\lesssim
\Xi_{\ell;\kappa_{\ell,\rho}\xi^2}
[\hsemig{\tau_1-\tau_0}{v_{\tau_0}}(X)]^2
\frac{L(L(1/\rho))}{L(1/\rho)}.
\end{equation}

Third, use the regularity of \(J_{\sigma,\rho}\) and the Gaussian-average regularity from Section 4 to replace the scale \(\Srho(\kappa_{\ell,\rho}\xi^2)\) and the flattened datum by the desired scale and datum. This contributes
\begin{equation}\label{eq:ch5-E3final}
E_3\lesssim
\left(
\Xi_{2;L(1/\rho)\xi^2}[\hsemig{\kappa_{\ell,\rho}^{-1}\xi^2}{v_{\tau_0}}(X)]
+\Xi_{2;t}[v_0]
\right)^2
\frac{L(L(1/\rho))}{L(1/\rho)}.
\end{equation}
Finally Proposition 4.4 controls all moment factors uniformly under the condition
\(\overline\beta>\eta(\beta,\ell)\). Adding \(E_1,E_2,E_3\) gives \eqref{eq:ch5-theorem514}.
\end{proof}

\subsection{The approximate root decoupling function approximates the root decoupling function}

For \(T_0<T_1\), define the exponential time change
\begin{align}\label{eq:ch5-RU}
R_{T_0,T_1,\rho}(q)
&:=T_0+(1-e^{-qL(1/\rho)})(T_1+\rho-T_0),\\
U_{T_0,T_1,\rho}(t)
&:=\frac1{L(1/\rho)}
\log\frac{T_1-T_0+\rho}{T_1-t+\rho}.
\end{align}
They are inverse maps on the relevant interval.

\begin{proposition}[Convergence of the approximate decoupling function]\label{prop:ch5-Jconv}
Let \(M,\lambda>0\), let \(\sigma\in\Lambda(M,\lambda)\), and fix
\(Q_0\in(0,\lambda^{-2})\). Then there exists
\(C=C(M,\lambda,Q_0)\) such that
\begin{equation}\label{eq:ch5-Jconv}
\sup_{Q\in[0,Q_0]}
\norm{(J_{\sigma,\rho}-J_{\sigma})(Q,\cdot)}_{\mathcal X}
\le
C\left[
\frac{L(L(1/\rho))}{L(1/\rho)}
\right]^{1/2}.
\end{equation}
\end{proposition}

\begin{proof}
Fix \(Q\in[0,Q_0]\) and constant initial data \(v_0\equiv a\). Choose \(T_0,T_1\) so that
\[
T_1-T_0=\Trho(Q),
\]
with the additional buffer needed to apply Theorem~\ref{thm:ch5-main-local}. The process
\[
V_t^{\rho,T_1}(x)=\hsemig{T_1-t}{v_t}(x)
\]
is a martingale. After replacing its quadratic variation by the box average of \(\sigma^2\), Theorem~\ref{thm:ch5-main-local} shows that the resulting coefficient is close to
\[
J_{\sigma,\rho}(Q-q,V_{R_{T_0,T_1,\rho}(q)}^{\rho,T_1}(x)).
\]
A white-noise coupling then gives an approximate SDE with this coefficient. In other words, \(J_{\sigma,\rho}\) is an approximate fixed point of the FBSDE map \(\mathcal Q_\sigma\):
\begin{equation}\label{eq:ch5-approx-fixedpoint}
\abs{J_{\sigma,\rho}(Q,a)-\mathcal Q_\sigma J_{\sigma,\rho}(Q,a)}_F^2
\le
C\langle a\rangle^2
\frac{L(L(1/\rho))}{L(1/\rho)}.
\end{equation}
The quantitative fixed-point stability estimate from Chapter 2 then yields
\[
\sup_{Q\le Q_0}
\frac{|J_{\sigma,\rho}(Q,a)-J_\sigma(Q,a)|_F}{\langle a\rangle}
\lesssim
\left[\frac{L(L(1/\rho))}{L(1/\rho)}\right]^{1/2},
\]
which is \eqref{eq:ch5-Jconv}.
\end{proof}

\subsection{Summary of the local mechanism}

The logical flow of Chapter 5 is
\[
\boxed{
\begin{array}{c}
\text{turn off the noise}\\
+\\
\text{flatten after a quiet interval}
\end{array}}
\quad\Longrightarrow\quad
\boxed{\text{localization in space}}
\quad\Longrightarrow\quad
\boxed{\text{concentration of }\sigma^2}
\]
\[
\Longrightarrow\quad
\boxed{
(S_{\xi,z}[\sigma^2\circ v_t])^{1/2}
\approx J_{\sigma,\rho}(q,\text{coarse field})}
\quad\Longrightarrow\quad
\boxed{J_{\sigma,\rho}\approx J_\sigma}.
\]
The restriction \(q<\lambda^{-2}\) is the genuinely local part of the argument. Chapter 6 removes this restriction by renormalization and induction on scales.

\begin{thebibliography}{9}
\bibitem{dunlap2026renormalizationflow2dnonlinear}
A.~Dunlap and C.~Graham,
\emph{Renormalization flow for the 2D nonlinear stochastic heat equation: pointwise statistics and universality},
2026, arXiv:2308.11850v3.
\end{thebibliography}

\end{document}
